\documentclass{article}
\usepackage{amsmath, amssymb, amsthm}
\usepackage{hyperref}
\usepackage{geometry}
\geometry{margin=1in}

\title{Erd\H{o}s Problem \#425}
\date{Last updated: 2025-08-31}
\author{Source: erdosproblems.com}

\begin{document}
\maketitle

\noindent\textbf{Status:} OPEN \quad \textbf{No prize}

\noindent\textbf{Tags:} number theory, sidon sets

\medskip

\noindent\textbf{Problem Statement:}

Let $F(n)$ be the maximum possible size of a subset $A\subseteq\{1,\ldots,N\}$ such that the products $ab$ are distinct for all $a&#60;b$. Is there a constant $c$ such that\[F(n)=\pi(n)+(c+o(1))n^{3/4}(\log n)^{-3/2}?\]If $A\subseteq \{1,\ldots,n\}$ is such that all products $a_1\cdots a_r$ are distinct for $a_1&#60;\cdots &#60;a_r$ then is it true that\[\lvert A\rvert \leq \pi(n)+O(n^{\frac{r+1}{2r}})?\]




Erd\H{o}s \cite{Er68} proved that there exist some constants $0&lt;c_1\leq c_2$ such that\[\pi(n)+c_1 n^{3/4}(\log n)^{-3/2}\leq F(n)\leq \pi(n)+c_2 n^{3/4}(\log n)^{-3/2}.\]This problem can also be considered in the real numbers: that is, what is the size of the the largest $A\subset [1,x]$ such that for any distinct $a,b,c,d\in A$ we have $\lvert ab-cd\rvert \geq 1$? Erd\H{o}s had conjectured (see \cite{Er73} and \cite{Er77c}) that $\lvert A\rvert=o(x)$. In \cite{ErGr80} Erd\H{o}s and Graham report that Alexander had given a construction disproving this conjecture, establishing that $\lvert A\rvert\gg x$ is possible. Alexander's construction is given in \cite{Er80}, and we sketch a simplified version. Let $B\subseteq [1,X^2]$ be a Sidon set of integers of size $\gg X$ and\[A=\{ X e^{b/X^2} : b\in B\}.\]It is easy to check that\[\lvert ab-cd\rvert \geq X^2\lvert 1-e^{1/X^2}\rvert \gg 1\]for distinct $a,b,c,d\in A$, and (after rescaling $A$ by some constant factor) this produces a set of size $\gg X$ with the desired property in the interval $[X,O(X)]$. Furthermore, a simple modification allows for $A$ to also be $1$-separated. In \cite{Er77c} Erd\H{o}s considers a similar generalisation for sets of complex numbers or complex integers. See also [490] , [793] , and [796] .




References

[Er68] Erd\H{o}s, P., On some applications of graph theory to number theoretic problems . Publ. Ramanujan Inst. (1968/69), 131-136.

[Er73] Erd\H{o}s, P., Problems and results on combinatorial number theory . A survey of combinatorial theory (Proc. Internat. Sympos., Colorado State Univ., Fort Collins, Colo., 1971) (1973), 117-138.

[Er77c] Erd\H{o}s, Paul, Problems and results on combinatorial number theory. III . Number theory day (Proc. Conf., Rockefeller Univ., New York, 1976) (1977), 43-72.

[Er80] Erd\H{o}s, Paul, A survey of problems in combinatorial number theory . Ann. Discrete Math. (1980), 89-115.

[ErGr80] Erd\H{o}s, P. and Graham, R., Old and new problems and results in combinatorial number theory . Monographies de L'Enseignement Mathematique (1980).

\noindent\textbf{Related OEIS sequences:} possible


\bigskip
\noindent\small{Source: \url{https://www.erdosproblems.com/425}}
\end{document}
